\section{Results}

\subsection{Visual Summary}
The first result to present is visual rather than only numerical: compare the
reconstructed diffraction modulus, amplitude, phase, and support for the same
sample IDs across model variants. This directly shows whether a lower
far-field metric corresponds to a plausible object-domain reconstruction.

\placeholderfigure{fig_autophasenn_reproduction.png}{AutoPhaseNN reproduction visualization.}

\placeholderfigure{fig_umamba_original.png}{UMamba original hard-threshold visualization.}

\placeholderfigure{fig_umamba_soft_threshold.png}{UMamba soft-threshold visualization.}

\placeholderfigure{fig_training_curves.png}{Training and validation loss curves from TensorBoard event files.}

\begin{table}[!t]
  \caption{Evaluation Metrics of Model Variants}
  \label{tab:main_results}
  \centering
  \begin{tabular}{lcccc}
    \toprule
    Model & FT L1 $\downarrow$ & Amp L1 $\downarrow$ & Phase L1 $\downarrow$ & Support IoU $\uparrow$ \\
    \midrule
    AutoPhaseNN & \todo{} & \todo{} & \todo{} & \todo{} \\
    UMamba & \todo{} & \todo{} & \todo{} & \todo{} \\
    Diagnostic ablation & \todo{} & \todo{} & \todo{} & \todo{} \\
    \bottomrule
  \end{tabular}
\end{table}

\subsection{AutoPhaseNN Reproduction}
The AutoPhaseNN reproduction is the anchor experiment. The expected narrative is:
the reproduced model should learn a compact object-domain amplitude/support
pattern and a phase field that remains spatially tied to the object region. The
paper should report both the far-field loss and the object-domain metrics, then
use the visualization in Fig.~\ref{fig:fig_autophasenn_reproduction.png} as the
qualitative baseline.

\todo{Insert final AutoPhaseNN checkpoint path, training epoch count, FT L1,
amplitude metrics, phase metrics, support metrics, and figure source.}

\subsection{UMamba Replacement}
The UMamba replacement is evaluated under the same data and physical
post-processing pipeline. The key observation to document is whether the
far-field metric improves without a corresponding improvement in amplitude,
phase, and support. If so, the result should be framed as an under-constrained
phase retrieval solution rather than simply as a failed optimizer run.

\todo{Insert UMamba checkpoint path, training epoch count, FT metrics, real-space
metrics, and representative visualization.}

\subsection{Diagnostic Ablation}
Only one ablation should be kept in the main paper: the one that most clearly
supports the explanation of the UMamba failure mode. For example, if the
soft-threshold experiment obtains acceptable numerical metrics but visibly
incorrect amplitude and phase, it can be used to show that smoother support
gradients alone do not solve the missing-prior problem. If a center-support
constraint is more informative, use that instead.

\todo{Choose one diagnostic ablation and report the same metric groups as the
other rows. Do not include every attempted training variant.}
